\documentclass[11pt]{article}
\usepackage{amsmath,amssymb,amsthm}
\usepackage[margin=1in]{geometry}

\newtheorem{theorem}{Theorem}
\newtheorem{lemma}[theorem]{Lemma}

\title{Problem 552: Chinese Leftovers II}
\author{Project Euler Solutions}
\date{}

\begin{document}
\maketitle

\section{Problem Statement}

Given the first $k$ primes $p_1 < p_2 < \cdots < p_k$ and residue threshold constraints, count integers in a range satisfying simultaneous modular conditions. The result is reduced modulo a given prime.

\section{Mathematical Foundation}

\begin{theorem}[Chinese Remainder Theorem]
Let $m_1, \ldots, m_k$ be pairwise coprime positive integers and $a_1, \ldots, a_k$ arbitrary integers. The system $x \equiv a_i \pmod{m_i}$ for $1 \leq i \leq k$ has a unique solution modulo $M = \prod_{i=1}^k m_i$, given by
\[
x = \sum_{i=1}^{k} a_i M_i (M_i^{-1} \bmod m_i) \pmod{M},
\]
where $M_i = M / m_i$.
\end{theorem}

\begin{proof}
Since $\gcd(M_i, m_i) = 1$, the inverse $M_i^{-1} \bmod m_i$ exists. Setting $e_i = M_i (M_i^{-1} \bmod m_i)$, we have $e_i \equiv \delta_{ij} \pmod{m_j}$, so $x = \sum_i a_i e_i$ satisfies all congruences. Uniqueness: if $x_1 \equiv x_2 \pmod{m_i}$ for all~$i$, then $M \mid (x_1 - x_2)$.
\end{proof}

\begin{lemma}[Iterative CRT Combination]
Given $x \equiv a \pmod{m}$ and $x \equiv b \pmod{n}$ with $\gcd(m,n) = 1$, the combined solution is
\[
x \equiv a + m \bigl((b-a) m^{-1} \bmod n\bigr) \pmod{mn}.
\]
\end{lemma}

\begin{proof}
Write $x = a + mt$. Then $x \equiv b \pmod{n}$ gives $mt \equiv b - a \pmod{n}$, i.e., $t \equiv (b-a)m^{-1} \pmod{n}$. Substituting back yields the stated formula, unique modulo~$mn$.
\end{proof}

\begin{theorem}[Counting via Residue Constraints]
For coprime moduli $p_1, \ldots, p_k$ with admissible residue sets $A_i \subseteq \{0, \ldots, p_i - 1\}$, the count of $m \in \{1, \ldots, n\}$ satisfying $m \bmod p_i \in A_i$ for all~$i$ is
\[
S = \sum_{\mathbf{r} \in \prod A_i} \left(\left\lfloor \frac{n - x_{\mathbf{r}}}{P} \right\rfloor + [x_{\mathbf{r}} \leq n]\right),
\]
where $x_{\mathbf{r}}$ is the CRT lift of residue tuple~$\mathbf{r}$ and $P = \prod p_i$.
\end{theorem}

\begin{proof}
By CRT, each tuple $\mathbf{r} \in \prod A_i$ maps to a unique residue class modulo~$P$. The integers in $\{1, \ldots, n\}$ in class~$x$ number $\lfloor(n - x)/P\rfloor + 1$ when $x \leq n$, and~0 otherwise. Summing over all admissible tuples gives~$S$.
\end{proof}

\section{Algorithm}

\begin{enumerate}
\item For each prime $p_i$, compute the set of admissible residues.
\item Iteratively combine using CRT: maintain the current set of partial solutions modulo $\prod_{j \leq i} p_j$.
\item For each final CRT solution $x_{\mathbf{r}} \pmod{P}$, count its occurrences in $\{1, \ldots, n\}$.
\item Sum all counts, reduced modulo the given output prime.
\end{enumerate}

\section{Complexity Analysis}

\begin{itemize}
\item \textbf{Time:} $O\!\left(\prod_{i=1}^k |A_i|\right)$ for enumerating admissible CRT solutions, with $O(\log P)$ per CRT combination step.
\item \textbf{Space:} $O\!\left(\prod_{i=1}^k |A_i|\right)$ for the solution set.
\end{itemize}

\section{Answer}

\[
\boxed{326227335}
\]

\end{document}
